
\documentclass[12pt]{article}
\pagestyle{empty} \topmargin -0.5in \textheight 9in \textwidth 6.5in
\oddsidemargin -0in
\usepackage{amssymb}
\usepackage{amsmath}
\usepackage[normalem]{ulem}
\usepackage{enumerate}
\usepackage{xcolor}

\DeclareMathOperator{\tr}{tr}

\newcommand{\vect}[1]{\mathbf{#1}}

\begin{document}

\begin{flushleft}

\textbf{Department of Applied Mathematics and Statistics}

\textbf{COLORADO SCHOOL OF MINES}\\


MATH $455$: Partial Differential Equations
\end{flushleft}

\vspace{0.1in}

\begin{center}
{\bf Assignment \#4 - Heat Equation} \\
Due Tuesday, October 20, 2026 \\
\end{center}

\vspace{0.2in}


\noindent $\textbf{1}$. (10 points) Let $k, L > 0$ and $f \in \mathrm{PC}^1[0,L]$ be a given function. Consider the heat equation
\begin{equation}
\label{HE}
\tag{HE}
\frac{\partial^2 u}{\partial x^2} - \frac{1}{k}\frac{\partial u}{\partial t}= 0.
\end{equation}
%
\begin{enumerate}[(a)]
\item Write down the solution to \eqref{HE} with the boundary conditions
$$\frac{\partial u}{\partial x}(t,0) = \frac{\partial u}{\partial x}(t,L) =0 , \qquad t > 0$$
and the initial condition
$$ u(0,x) = f(x), \qquad 0 < x < L.$$ \\
\textcolor{blue}{\textbf{Solution} — As with the wave equation, we apply an identical strategy of assuming the solution takes the separated form $$u(t,x) = T(t)X(x)$$ Then $$\frac{\partial^2 u}{\partial x^2} = T(t)X''(x) \;\;\;\;\;\;\;\; \frac{\partial u}{\partial t} = T'(t)X(x)$$ \\ and plugging into the heat equation, we obtain: \\ $$k\frac{X''(x)}{X(x)} = \frac{T'(t)}{T(t)} = -\lambda$$ The left-hand side is a function of $x$ only and the right-hand side is a function of $t$ only. For them to be equal to one another for all values of $t,x$, they must both be equal to the same constant, which (for future convenience) we will notate as $-\lambda$. Hence our separated heat equation reduces to a system of independent, linear, homogeneous ordinary differential equations:
$$\begin{cases}
kX''(x) + \lambda X(x) = 0 \\
T'(t) + \lambda T = 0
\end{cases}$$\\
We see that the structure of solutions to these equations depends on the sign of $\lambda$, so we split into cases accordingly: \\ \\
\textbf{For $\lambda < 0$:} \\
$$T(t) = c_0 e^{-\lambda t} \;\;\;\;\;\;\; X(x) = c_1 e^{\sqrt{-\frac{\lambda}{k}} x} + c_2 e^{-\sqrt{-\frac{\lambda}{k}} x}$$
\textbf{For $\lambda = 0$:} \\
$$T(t) = c_3 \;\;\;\;\;\;\;  X(x) = c_4 + c_5x$$
\textbf{For $\lambda > 0$:} \\ 
$$T(t) = c_6 e^{-\lambda t} \;\;\;\;\;\;\; X(x) = c_7 \cos \left( \sqrt{\frac{\lambda}{k}} x \right) + c_8 \sin \left( \sqrt{\frac{\lambda}{k}} x \right)$$ \\ \\
We now constrain these general solutions with the homogenous \textbf{Neumann} boundary conditions: \\ $$\frac{\partial u}{\partial x} (t,0) = 0 \implies T(t)X'(0) = 0 \implies X'(0) = 0 $$ 
$$\frac{\partial u}{\partial x}(t,L) = 0 \implies T(t)X'(L) = 0 \implies X'(L) = 0$$ \\
\textbf{For $\lambda < 0$}:
\begin{align*}
X'(0) = 0 \implies& c_1 \sqrt{-\frac{\lambda}{k}} - c_2 \sqrt{-\frac{\lambda}{k}} = 0 \implies \boxed{c_1 = c_2} \\
X'(L) = 0 \implies& c_1 \sqrt{-\frac{\lambda}{k}} e^{\sqrt{-\frac{\lambda}{k}} L} - c_2 \sqrt{-\frac{\lambda}{k}} e^{-\sqrt{-\frac{\lambda}{k}} L} = 0 \\ 
\implies& c_1 \sqrt{-\frac{\lambda}{k}} \left( e^{\sqrt{-\frac{\lambda}{k}} L} - e^{-\sqrt{-\frac{\lambda}{k}} L}   \right) = 0 \;\;\;\; \text{(since $c_1 = c_2$)} \\
\implies& \boxed{c_1 = 0} \;\;\; \text{(Exponential factor $\neq 0$ for $\lambda \neq 0$)}
\end{align*}
\textbf{For $\lambda = 0$:} \\
\begin{align*}
X'(0) = X'(L) = 0 \implies& \boxed{c_5 = 0} \\
\end{align*}
\textbf{For $\lambda > 0$:} \\
\begin{align*}
X'(0) = 0 \implies& c_8 \sqrt{\frac{\lambda}{k}} = 0 \\
\implies& \boxed{c_8 = 0} \;\;\; \text{(since $\lambda > 0, k > 0$)} \\
X'(L) = 0 \implies& -c_7 \sqrt{\frac{\lambda}{k}} \sin \left( \sqrt{\frac{\lambda}{k}} L \right) = 0 \\
\implies& \sin \left( \sqrt{\frac{\lambda}{k}} L \right) = 0 \;\;\; \text{(since $\lambda \neq 0, k \neq 0$, and assuming $c_7 \neq 0$)} \\
\end{align*}
Note that $c_7 = 0$ would produce a trivial solution since $c_8$ already equals zero, and the trivial solution is already accounted for in the other $\lambda$-cases.
\begin{align*}
\implies& \sqrt{\frac{\lambda}{k}} L = n\pi \;\;\; \text{(for any $n \in \mathbb{Z}$)} \\
\implies& \boxed{\lambda = k \left( \frac{n\pi}{L} \right)^2}
\end{align*}
We write the general solutions produced by each case (with the Neumann boundary conditions applied): \\ \\
\textbf{For $\lambda < 0$}: \\ \\
$$c_1 = c_2 = 0 \implies X(x) = 0 \implies u(t,x) = 0$$
\textbf{For $\lambda = 0$}: \\ \\
$$c_5 = 0 \implies X(x) = c_4 \implies u(t,x) = c_3 c_4 = C_0 \;\;\; \text{for some $C_0 \in \mathbb{R}$}$$ \\ \\
\textbf{For $\lambda > 0$}: \\ \\
$$c_8 = 0 \implies u_n(t,x) = C_n e^{- k \left( \frac{n\pi}{L} \right)^2 t} \cos \left( \frac{n\pi x}{L} \right) \;\;\; \text{for $C_n \in \mathbb{R}$}$$ \\
Note that each separated eigenfunction solution will have its own constant $C_n$. Hence the general solution is the superposition of all solutions from each of these 3 cases:
$$\boxed{u(t,x) = C_0 + \sum_{n=1}^{\infty} C_n e^{- k \left( \frac{n\pi}{L} \right)^2 t} \cos \left( \frac{n\pi x}{L} \right)}$$ \\ \\
Let us now apply the initial condition $u(0,x) = f(x)$ for $0 \leq x \leq L$.
$$\implies C_0 + \sum_{n=1}^{\infty} C_n \cos \left( \frac{n\pi x}{L} \right) = f(x)$$
Notice that the application of this initial condition creates an equation that is exactly the Fourier series expression of $f(x)$. For both sides to be equal, we must have that each $C_n$ is exactly the $n$th Fourier coefficient in the expansion of $f(x)$ on $[0,L]$. Hence $$\boxed{C_0 = \frac{1}{L} \int_{0}^L f(x) dx}$$
$$\boxed{C_n = \frac{2}{L} \int_{0}^L f(x) \cos \left( \frac{n \pi x}{L} \right) dx \;\;\;\;\; n \geq 1}$$
}
\item Let $T_1 > 0$, take
$$f(x) = T_1 \frac{x}{L},$$
and determine the resulting solution by computing the Fourier coefficients and the corresponding Fourier series. \\ \\ \textcolor{blue}{\textbf{Solution} — 
\begin{align*}
C_0 &= \frac{1}{L} \int_0^L \left( T_1 \frac{x}{L} \right) dx \\
&= \frac{T_1}{L^2} \int_0^L x dx \\
&= \frac{T_1}{L^2} \left( \frac{1}{2} L^2 \right) \\
&= \frac{T_1}{2}
\end{align*}
For the general coefficients
\begin{align*}
C_n &= \frac{2}{L} \int_0^L \left( T_1 \frac{x}{L} \right) \cos \left( \frac{n\pi x}{L} \right) dx \\
&= \frac{2T_1}{L^2} \int_0^L x \cos \left( \frac{n\pi x}{L} \right) dx
\end{align*}
Apply integration by parts:
\begin{align*}
U = x \;\;\;\;\;\;\;\;& dV = \cos \left( \frac{n\pi x}{L} \right) dx \\
dU = dx \;\;\;\;\;\;\;\;& V = \frac{L}{n\pi} \sin \left( \frac{n\pi x}{L} \right)
\end{align*}
\begin{align*}
C_n &= \frac{2T_1}{L^2} \left[ \underbrace{ \left. \frac{Lx}{n\pi} \sin \left( \frac{n\pi x}{L} \right) \right|_{x=0}^{x=L}}_{0} -  \frac{L}{n\pi} \int_0^L \sin \left( \frac{n\pi x}{L} \right) dx   \right] \\
&= \frac{-2T_1}{n\pi L} \left[ - \frac{L}{n\pi} \cos \left( \frac{n \pi x}{L} \right) \right]_{x=0}^{x=L} \\ 
&= \frac{2T_1}{(n\pi)^2} \left[ \cos \left( \frac{n \pi x}{L} \right) \right]_{x=0}^{x=L} \\
&= \frac{2T_1}{(n\pi)^2} \left[ \cos (n\pi) - 1 \right] \\
&= \frac{2T_1}{(n\pi)^2} \left[ (-1)^n - 1 \right] \\
\end{align*}
Hence the solution is $$\boxed{u(t,x) = \frac{T_1}{2} + \frac{2T_1}{\pi^2} \sum_{n=1}^{\infty} \left[ \frac{(-1)^n - 1}{n^2} \right] e^{- k \left( \frac{n\pi}{L} \right)^2 t} \cos \left( \frac{n\pi x}{L} \right)}$$
}
\end{enumerate}

\vspace{0.3in}


\noindent $\textbf{2}$. (15 points)  Let $L, T_0 > 0$ and $f \in \mathrm{PC}^1[0,L]$ be given. Consider the temperature distribution of a long, thin bar in which heat is internally generated by a source with constant rate.
In particular, we let $h > 0$ be given and consider the heat equation
$$\frac{1}{k}\frac{\partial u}{\partial t} = \frac{\partial^2 u}{\partial x^2} + h$$ 
for $0 < x < L$ and $t > 0$ 
with given boundary conditions
$$u(t, 0) = T_0 \qquad \mathrm{and} \qquad \frac{\partial u}{\partial x}(t, L) = 0$$
and initial temperature distribution
$$u(0,x) = f(x).$$
Determine the temperature throughout the bar at all later times, namely compute $u(t,x)$. \\ \\
\textcolor{blue}{\textbf{Solution} — Because this is the inhomogeneous heat equation (caused by the presence of the $h$ term), we cannot separate the equation that results from assuming $u(t,x) = T(t)X(x)$ and plugging in. Hence we need a different method—namely, the one that involves finding a steady-state solution. \\ \\ Assume the bar settles on the steady-state distribution $u_s(x)$ as $t \to \infty$. Because this particular variation of the heat equation should be valid for all time, it's definitely valid for $u(t,x) = u_s(x)$ at $t = \infty$. Hence, plugging in $u_s(x)$, we obtain
\begin{align*}
&0 = u_s''(x) + h \\
\implies& u_s''(x) = -h\\
\implies& u_s'(x) = -hx + C_0\\
\implies& u_s(x) = -\frac{1}{2} hx^2 + C_0 x + C_1\\
\end{align*}
We can apply the boundary conditions to the steady-state solution:
\begin{align*}
u(t,0) = T_0 \implies& u_s(0) = T_0 \\
\implies& C_1 = T_0 \\
\frac{\partial u}{\partial x}(t, L) = 0 \implies& u_s'(L) = 0 \\
\implies& -hL + C_0 = 0 \\
\implies& C_0 = hL \\
\end{align*}
Hence the steady-state solution with applied boundary conditions is $$u_s(x) = -\frac{1}{2} h x^2 + hLx + T_0$$ We can now perform a change-of-variable transformation that uses this steady-state solution to transform the given heat equation into a homogeneous (and separable) form. Namely, let $v(t,x)$ be the difference in heat between the current heat $u(t,x)$ and the eventual final heat $u_s(x)$: $$v(t,x) = u(t,x) - u_s(x)$$
We plug this new expression of $u(t,x)$ into the given heat equation (Note that $u = v + u_s$)
$$\frac{1}{k} \frac{\partial v}{\partial t} = \frac{\partial^2 v}{\partial x^2} + u_s''(x) + h$$
And note that $u_s''(x) = -h$, so we recover the homogeneous form
$$\frac{1}{k} \frac{\partial v}{\partial t} = \frac{\partial^2 v}{\partial x^2}$$
With boundary conditions 
\begin{align*}
u(t,0) = T_0 \implies& v(t,0) + \underbrace{u_s(0)}_{T_0} = T_0 \\
\implies& v(t,0) = 0 \\
\end{align*}
\begin{align*}
\frac{\partial u}{\partial x} (t,L) = 0 \implies& \frac{\partial v}{\partial x}(t,L) + \underbrace{u_s'(L)}_{0} = 0\\
\implies& \frac{\partial v}{\partial x}(t,L) = 0
\end{align*}
And initial condition
\begin{align*}
u(0,x) = f(x) \implies& v(0,x) + u_s(x) = f(x) \\
\implies& v(0,x) = f(x) - u_s(x) \\
\end{align*}
Since we've transformed the problem to a homogeneous form in $v(t,x)$, we proceed to solve regularly in $v$, assuming $v(t,x) = T(t)X(x)$. Then $$k\frac{X''(x)}{X(x)} = \frac{T'(t)}{T(t)} = -\lambda$$ \\
$$\begin{cases}
kX''(x) + \lambda X(x) = 0 \\
T'(t) + \lambda T = 0
\end{cases}$$\\
The structures of solutions here are, again, dependent on the sign of $\lambda$, so we examine each one: \\ \\
\textbf{For $\lambda < 0$:} \\
$$T(t) = c_0 e^{-\lambda t} \;\;\;\;\;\;\; X(x) = c_1 e^{\sqrt{-\frac{\lambda}{k}} x} + c_2 e^{-\sqrt{-\frac{\lambda}{k}} x}$$
\textbf{For $\lambda = 0$:} \\
$$T(t) = c_3 \;\;\;\;\;\;\;  X(x) = c_4 + c_5x$$
\textbf{For $\lambda > 0$:} \\ 
$$T(t) = c_6 e^{-\lambda t} \;\;\;\;\;\;\; X(x) = c_7 \cos \left( \sqrt{\frac{\lambda}{k}} x \right) + c_8 \sin \left( \sqrt{\frac{\lambda}{k}} x \right)$$ \\ \\
We apply the boundary conditions: \\ \\
\textbf{For $\lambda < 0$}: \\ \\
\begin{align*}
v(t,0)& = 0 \implies T(t)X(0) = 0  \\
\implies& X(0) = 0 \;\;\; \text{(Assuming $T(t) \neq 0$)} \\ 
\implies& \boxed{c_1 = -c_2} \\ \\
\frac{\partial v}{\partial x}(t,L) & = 0 \\
\implies& T(t)X'(L) = 0 \\
\implies& X'(L) = 0 \;\;\; \text{(Assuming $T(t) \neq 0$)} \\
\implies& c_1 \sqrt{-\frac{\lambda}{k}} e^{\sqrt{-\frac{\lambda}{k}} L} - c_2 \sqrt{-\frac{\lambda}{k}} e^{-\sqrt{-\frac{\lambda}{k}} L} = 0 \\ 
\implies&  c_1 \sqrt{-\frac{\lambda}{k}} \left( \underbrace{ e^{\sqrt{-\frac{\lambda}{k}} L} + e^{-\sqrt{-\frac{\lambda}{k}} L}}_{\neq 0}  \right) \\
\implies& \boxed{c_1 = 0} \;\;\; \text{(since $\lambda \neq 0, k \neq 0$)}
\end{align*}
\textbf{For $\lambda = 0$}: \\ \\
\begin{align*}
X(0) = 0 \implies& \boxed{c_4 = 0} \\
X'(L) = 0 \implies& \boxed{c_5 = 0}
\end{align*}
\textbf{For $\lambda > 0$}: \\ \\
\begin{align*}
X(0) = 0 \implies& \boxed{c_7 = 0} \\
X'(L) = 0 \implies& c_8 \sqrt{\frac{\lambda}{k}} \cos \left( \sqrt{\frac{\lambda}{k}} L \right) = 0 \\
\implies& \sqrt{\frac{\lambda}{k}} L = \frac{\pi}{2} + n\pi \;\;\; \text{for $n \in \mathbb{Z}$ (assuming $c_8 \neq 0$)} \\
\implies& \lambda = k \left( \frac{\pi}{2L} + \frac{n\pi}{L} \right)^2 \\ 
\implies& \lambda = \frac{k \pi^2}{L^2} \left( \frac{1}{2} + n \right)^2 \\ 
\implies& \lambda = \frac{k \pi^2}{4L^2} \left(2n + 1 \right)^2 \\ 
\end{align*}
Our boundary condition-constrained constants cause $X(x) = 0$ and hence $u(t,x) = 0$ for $\lambda < 0$ and $\lambda = 0$. Hence, the superposition of all separated solutions consists only of $\lambda > 0$: \\ \\ 
$$v(t,x) = \sum_{n=0}^{\infty} C_n e^{- \frac{k \pi^2  \left(2n + 1 \right)^2}{4L^2} t} \sin \left( \frac{(2n+1)\pi x}{2L} \right)$$ \\
We transform back to $u(t,x)$, using $u(t,x) = u_s(x) + v(t,x)$: \\
$$u(t,x) = -\frac{1}{2} h x^2 + hLx + T_0 + \sum_{n=0}^{\infty} C_n e^{- \frac{k \pi^2  \left(2n + 1 \right)^2}{4L^2} t} \sin \left( \frac{(2n+1)\pi x}{2L} \right)$$ \\
We now apply the initial condition \\
\begin{align*}
u(0,x) = f(x) \implies& -\frac{1}{2} h x^2 + hLx + T_0 + \sum_{n=0}^{\infty} C_n \sin \left( \frac{(2n+1)\pi x}{2L} \right) = f(x)\\
\implies& \sum_{n=0}^{\infty} C_n \sin \left( \frac{(2n+1)\pi x}{2L} \right) = f(x) + \frac{1}{2} h x^2 - hLx - T_0 \\
\end{align*}
We can't read off the $C_n$'s as the coefficients of any identifiable Fourier series yet, because the summand isn't identical to the familiar $n\pi x / L$ frequency. Hence, we need to isolate the $C_n$'s manually using the orthogonality of the $\sin$ eigenfunctions. \\ \\ We multiply both sides by an equivalent $m$-frequency $\sin$ function and integrate from $0$ to $L$. We'll start with the left-hand side:
\begin{align*}
& \int_0^L \sin \left( \frac{(2m+1)\pi x}{2L}  \right) \sum_{n=0}^{\infty} C_n \sin \left( \frac{(2n+1)\pi x}{2L} \right) dx \\
&= \sum_{n=0}^{\infty}  C_n \int_0^L  \sin \left( \frac{(2m+1)\pi x}{2L}  \right) \sin \left( \frac{(2n+1)\pi x}{2L} \right) dx \\
\end{align*}
Only the $n=m$ term in this sum is nonzero:
$$= C_m \int_0^L  \sin \left( \frac{(2m+1)\pi x}{2L} \right)^2 dx$$
Let $$u = \frac{(2m+1)\pi x}{2L} \;\;\;\;\; du = \frac{(2m+1)\pi}{2L} dx$$  \\
The integral becomes
\begin{align*}
& C_m \frac{2L}{(2m + 1)\pi} \int_0^{\frac{(2m+1)\pi}{2}} \sin^2(u)du \\
&= C_m \frac{L}{(2m + 1)\pi} \int_0^{\frac{(2m+1)\pi}{2}} 1 - \cos(2u) du \;\;\; \text{(applying $\sin$ power reduction)}  \\
&= C_m \frac{L}{(2m + 1)\pi} \left( \frac{(2m+1)\pi}{2} \right) - \frac{L}{(2m+1)\pi} \underbrace{\left[ \frac{1}{2} \sin(2u) \right]_{0}^{\frac{(2m+1)\pi}{2}}}_{0} \\
&= C_m \frac{L}{2}
\end{align*} \\
Hence, our $0$-to-$L$-integrated, $m$-frequency-multiplied equation is 
\begin{align*}
& C_m \frac{L}{2} = \int_0^L \sin \left( \frac{(2m+1)\pi x}{2L}  \right) \left( f(x) + \frac{1}{2} h x^2 - hLx - T_0 \right) dx \\
\implies& C_m = \frac{2}{L} \int_0^L \sin \left( \frac{(2m+1)\pi x}{2L}  \right) \left( f(x) + \frac{1}{2} h x^2 - hLx - T_0 \right) dx \\
\end{align*} \\
Hence the temperature of the bar at all later times is \\
$$\boxed{u(t,x) = -\frac{1}{2} h x^2 + hLx + T_0 + \sum_{n=0}^{\infty} C_n e^{- \frac{k \pi^2  \left(2n + 1 \right)^2}{4L^2} t} \sin \left( \frac{(2n+1)\pi x}{2L} \right)}$$
With $$\boxed{C_n = \frac{2}{L} \int_0^L \sin \left( \frac{(2n+1)\pi x}{2L}  \right) \left( f(x) + \frac{1}{2} h x^2 - hLx - T_0 \right) dx}$$
}
\vspace{0.3in}

\noindent $\textbf{3}$. (13 points) Let $c$ be a given number satisfying $\frac{\pi}{10} < c < \frac{9\pi}{10}$ and define
for every $0 < x < \pi$ the function 
$$h(x) =
\begin{cases}
	5, & \mathrm{if} \ c - \frac{\pi}{10} < x < c + \frac{\pi}{10} \\ 
	0, & \mathrm{else}.
\end{cases}$$
%
\begin{enumerate}[(a)]
\item Solve the heat equation with $k=1$, namely
$$\frac{\partial^2 u}{\partial x^2} - \frac{\partial u}{\partial t}= 0$$
with
$$u(t,0) = u(t,\pi) =0 , \qquad t > 0$$
and
$$ u(0,x) = h(x), \qquad 0 < x < \pi.$$
Use the trigonometric identity
$$ \cos(\alpha) - \cos(\beta) = -2\sin \left ( \frac{\alpha + \beta}{2} \right )\sin \left ( \frac{\alpha - \beta}{2} \right )$$
to simplify the form of the Fourier coefficients. \\ \\
\textcolor{blue}{\textbf{Solution} — Since we already solved the general solution in problem 1, we solve the general heat equation identically (as before) and plug in $k=1$ later. The separated general solutions $u(t,x) = T(t)X(x)$ to the homogenous heat equation takes the forms (as we've seen before in problem 1): \\  \\
\textbf{For $\lambda < 0$:} \\
$$T(t) = c_0 e^{-\lambda t} \;\;\;\;\;\;\; X(x) = c_1 e^{\sqrt{-\frac{\lambda}{k}} x} + c_2 e^{-\sqrt{-\frac{\lambda}{k}} x}$$
\textbf{For $\lambda = 0$:} \\
$$T(t) = c_3 \;\;\;\;\;\;\;  X(x) = c_4 + c_5x$$
\textbf{For $\lambda > 0$:} \\ 
$$T(t) = c_6 e^{-\lambda t} \;\;\;\;\;\;\; X(x) = c_7 \cos \left( \sqrt{\frac{\lambda}{k}} x \right) + c_8 \sin \left( \sqrt{\frac{\lambda}{k}} x \right)$$ \\ \\
We now apply the \textbf{Dirichlet} boundary conditions: \\
\begin{align*}
u(t,0) = 0 \implies& T(t)X(0) = 0 \\
\implies& X(0) = 0 \;\;\; \text{(for non-trivial $T(t) \neq 0$)} \\
u(t,\pi) = 0 \implies& T(t)X(\pi) = 0 \\
\implies& X(\pi) = 0 \;\;\; \text{(for non-trivial $T(t) \neq 0$)} \\
\end{align*}
\textbf{For $\lambda < 0$:} \\ \\
\begin{align*}
X(0) = 0 \implies& \boxed{c_1 = -c_2} \\
X(\pi) = 0 \implies& c_1 e^{\sqrt{-\frac{\lambda}{k}} \pi} + c_2 e^{-\sqrt{-\frac{\lambda}{k}} \pi} = 0 \\
\implies& c_1 \left(e^{\sqrt{-\frac{\lambda}{k}} \pi} - e^{-\sqrt{-\frac{\lambda}{k}} \pi} \right) = 0 \;\;\; \text{(using $c_1 = -c_2$)} \\
\implies& \boxed{c_1 = 0}
\end{align*}
(only the trivial solution exists in this case) \\ \\
\textbf{For $\lambda = 0$:} \\ \\
\begin{align*}
X(0) = 0 \implies& \boxed{c_4 = 0} \\
X(\pi) = 0 \implies& c_4 + \pi c_5 = 0 \\
\implies& \pi c_5 = 0 \\
\implies& \boxed{c_5 = 0}
\end{align*}
(only the trivial solution exists in this case) \\ \\
\textbf{For $\lambda > 0$:} \\ \\
\begin{align*}
X(0) = 0 \implies& \boxed{c_7 = 0}\\
X(\pi) = 0 \implies& c_8 \sin \left( \sqrt{\frac{\lambda}{k}} \pi \right) = 0\\
\implies& \sqrt{\frac{\lambda}{k}} \pi = n\pi \;\;\; \text{for $n \in \mathbb{Z}$ (and assuming $c_8 \neq 0$)} \\
\implies& \lambda = kn^2
\end{align*}
Hence the general solution (the superposition of all our boundary-constrained solutions over all $\lambda$) is 
\begin{align*}
u(t,x) &= \sum_{n=0}^{\infty} C_n e^{-kn^2 t} \sin \left( nx \right) \\
\text{or} \;\;\;u(t,x) &= \sum_{n=0}^{\infty} C_n e^{-n^2 t} \sin \left( nx \right) \;\;\; \text{(with $k=1$)} \\
\end{align*}
We now apply the initial condition \\
\begin{align*}
& u(0,x) = h(x) \;\;\; \text{(for $0 \leq x \leq \pi$)} \\
\implies& \sum_{n=1}^{\infty} C_n \sin \left( nx \right) = h(x) \\
\end{align*}
This is exactly the Fourier $\sin$ expansion of $h(x)$ on $[0, \pi]$. Hence \\ \\
\begin{align*}
C_n &= \frac{2}{\pi} \int_0^{\pi} h(x) \sin (nx) dx \\
\end{align*}
Note that
\begin{align*}
& \frac{\pi}{10} < c < \frac{9\pi}{10} \\
\implies& c - \frac{\pi}{10} > 0 \;\; \text{and} \;\; c + \frac{\pi}{10} < \pi
\end{align*}
So the piecewise interval is contained fully in $[0,\pi]$: $$\left[ c - \frac{\pi}{10}, c + \frac{\pi}{10}\right] \in [0,\pi]$$ Hence
\begin{align*}
C_n &= \frac{2}{\pi} \int_{c - \frac{\pi}{10}}^{c + \frac{\pi}{10}} h(x) \sin (nx) dx
\end{align*}
}
\item Notice that for $t$ sufficiently large, we have
$$u(t,x) \approx  B_1 e^{-t} \sin(x).$$
What choice of $c$ as defined above makes $B_1$ largest?
\item Using the idea in the previous part, answer the following question:
Suppose that you're forced to put a cup of hot coffee somewhere between two blocks of ice. Where exactly should you put it if you want it to stay as warm as possible?
\end{enumerate}


\vspace{0.3in}


\noindent $\textbf{4}$. (12 points)  Let $f, g \in \mathrm{PC}^1[0,\pi]$ be given and find the solution of the inhomogeneous wave equation with speed of propagation $c=1$, namely
$$\frac{\partial^2 u}{\partial t^2} - \frac{\partial^2 u}{\partial x^2} +1 = 0$$
for $0 < x < \pi$ and $t > 0$ 
subject to the given boundary conditions
$$u(t, 0) = 0 \qquad \mathrm{and} \qquad u(t,\pi) = 0$$
for $t > 0$ and initial conditions
$$u(0,x) = f(x) \qquad \mathrm{and} \qquad \frac{\partial u}{\partial t} (0,x) = g(x)$$
for $0 < x < \pi$.

\end{document}
